%! TeX program = lualatex
\documentclass{perfil}
\usepackage{fontspec}
\newcommand{\tabitem}{~~\llap{\textbullet}~~}
\setmainfont{Times New Roman}
\newcommand{\periodo}{marzo 2024 - agosto 2024}

\begin{document}

%Portada - Indice de contenidos - Datos Informativos A

%Portada
\begin{titlepage}
	\begin{center}
		\begin{tabular}{l  p{0.65\textwidth}  l}
			\multirow{2}{*}{
				\includegraphics[width=20mm]{resources/images/logo-uta-rojo.png}
			}                                                             &
			\bfseries\uppercase{\centering Universidad Técnica de Ambato} &
			\multirow{2}{*}{
				\includegraphics[width=25mm]{resources/images/logo-fisei-feo.jpeg}
			}                                                                                                                                                              \\
			                                                              & \bfseries\uppercase{\centering Facultad de Ingeniería en Sistemas, Electrónica e Industrial} & \\
		\end{tabular}
		\vspace{2cm}
	\end{center}
	{\textbf{MODALIDAD: }\par}
	\begin{center}
		Proyecto de Investigación \fbox{X} \\
		Artículo Académico \fbox{\phantom{X}}
	\end{center}
	{\textbf{AUTOR: }\par}
	\begin{center}
		Pazmiño Fiallos Jonathan José
	\end{center}
	{\textbf{CARRERA: }\par}
	\begin{center}
		Software
	\end{center}
	{\textbf{LÍNEA DE INVESTIGACIÓN: }\par}
	\begin{center}
		Tecnología de la información y sistemas de control
	\end{center}
	{\textbf{PERIODO: }\par}
	\begin{center}
		\periodo
	\end{center}
	{\textbf{LUGAR Y FECHA DE PRESENTACIÓN: }\par}
	\begin{center}
		% Ambato, 18 de enero de 2023
	\end{center}
	\vfill
\end{titlepage}


%Numeracion en romanos minuscula y comenzando desde el ``2''
\pagenumbering{roman}
\setcounter{page}{2}

%Titulo de la tabla de contendios centrado, bold y mayusculas
\tableofcontents
\pagebreak

\listoffigures
\pagebreak

\listoftables
\pagebreak

%Numeros arabicos despues de los prelims
%Datos informativos
\begin{center}
	\textbf{A. DATOS INFORMATIVOS}
\end{center}\par
\textbf{Modalidad: } Trabajo de Titulación\par
\textbf{Opción de titulación: } Proyecto de Investigación\par
\textbf{Autor: } Pazmiño Fiallos Jonathan José\par
\textbf{Carrera: } Software\par
\textbf{Area: } Gestión de Tecnologías de la Información\par
\textbf{Línea de investigación: } Tecnología de la Información y Sistemas de Control\par
\textbf{Periodo: } \periodo\par
\textbf{Lugar y fecha de presentación: } Ambato,  17 de Enero de 2024
\pagebreak


\renewcommand{\thepage}{\arabic{page}}% Roman page numbers
\begin{center} \textbf{B. CONTENIDOS} \end{center}
\section{Tema de investigación}
\begin{nohyphens}
	\uppercase{ Desarrollo de un sistema web y móvil que mejore la obtención de
		datos georreferenciados de delitos con el propósito de crear un modelo
		analítico de Business Intelligence (BI) orientado al
		seguimiento de las tendencias de siniestros en la ciudad de Ambato.}
\end{nohyphens}
\subsection{Planteamiento del problema}
En \cite{rummensMobilePhoneData2021}, desde el 2020 se ha notado un significativo incremento en la cantidad de
campos de estudio que llevan a cabo el proceso de toma de decisiones basadas en la información extraída de grandes
conjuntos de datos. Según \cite{ahishakiyePerformanceAnalysisBusiness2017}, para asegurar que el proceso de toma
de decisiones conduzca a los resultados óptimos, es fundamental el uso de Business Intelligence (BI), el cual
emplea un conjunto de herramientas que facilitan la adquisición, almacenamiento, análisis y entrega de la información
necesaria.

\bigbreak
En \cite{reidDevelopingModelPerceptions2020} se menciona que en el ámbito social, el incremento de la delincuencia,
el miedo y la inseguridad ciudadana respecto a la zona en la que residen son temas que afectan directamente la calidad
de vida de las personas. En este contexto, \cite{rummensMobilePhoneData2021} señala que el análisis de datos sobre la
delincuencia ha empezado a tomar una mayor relevancia, siendo estos datos usados de forma objetiva para fundamentar
políticas, así como estrategias policiales y operaciones tácticas con el fin de reducir y prevenir la delincuencia
mediante la toma de decisiones oportunas y precisas.

\bigbreak
La carencia de fuentes adecuadas de datos georreferenciados, con los cuales realizar análisis, dificulta la capacidad de
la ciudadanía y las fuerzas del orden para tomar precauciones necesarias \cite{tasnimNovelMultiModuleApproach2022}. Según
\cite{vieiraCrimePredictionPrevention2022}, los datos georreferenciados o espacio-temporales son un valioso recurso utilizado
por analistas criminalistas. Estos datos son la base para la identificación de patrones delictivos en áreas de alta concentración
de delitos a través del mapeo de zonas críticas o puntos calientes (hotspots). El autor en \cite{rummensMobilePhoneData2021}
menciona que el uso de software permite la adquisición de datos delictivos de forma más precisa y eficiente, obteniendo información
útil que permita identificar los puntos focales del crimen.

\bigbreak
% georreferenciados para el control de la delincuencia \cite{ariasCrimePunishmentLatin2023}.
En el estudio de \cite{unitednationsReferenceFrameworkSecurity2022}, según la Comisión Económica para América Latina y el Caribe (CEPAL),
en América Latina y el Caribe, la georreferenciación se posiciona como la clave para realizar análisis más profundos y detallados sobre la
delincuencia. Esta técnica además facilita la visualización de la distribución espacio-temporal de variables relacionadas con el crimen,
permitiendo su análisis posterior. En ese sentido, el autor de \cite{ariasCrimePunishmentLatin2023} menciona que desde 1990 se han establecido
paradigmas de entes policiales orientados a la comunidad, los cuales emplean estrategias contra la delincuencia haciendo uso de datos
georreferenciados.

\bigbreak
Según \cite{tayupantaBOLETINSEMESTRALHOMICIDIOS2023}, el boletín semestral de homicidios intencionados en Ecuador, presentado por la
Fundación Panamericana para el Desarrollo (PADF), menciona que en 2022, el país registró un total de 4603 homicidios intencionales,
marcando así el registro histórico más alto hasta el momento. Se observó un promedio diario de 10 casos, junto a una tasa de 25.9
homicidios por cada 100 000 habitantes. Sin embargo, en el primer semestre de 2023, los casos de homicidios intencionados registrados
descendieron a 3599, lo que equivale a un promedio diario de 19 casos, casi el doble de los registros obtenidos en 2022.
Debido a esta tendencia, se prevé que la cifra total de homicidios supere los 7000 casos, alcanzando una tasa de 35 homicidios por
cada 100 000 habitantes. Esto sitúa a Ecuador como uno de los tres países más violentos en América Latina, solo por detrás de Venezuela y
Honduras. Por lo tanto, en \cite{narvaezGESTIONINFORMACIONCRIMINOLOGICA2018} se plantea la necesidad de abordar activamente la seguridad
en beneficio de la ciudadanía ecuatoriana, proponiendo el uso de tecnologías que permitan establecer estadísticas, interfaces gráficas
y mapas con información georreferenciada.

% Desde el año 2011 el Departamento de Análisis de Información del Delito (DAID) de la policía nacional ha
% trabajado en la implementación de un sistema de información respecto a los incidentes ocurridos en el territorio nacional.


% El sistema
% David desarrollado por el DAID se encarga de generar información actualizada de crímenes la cual es usada como base para la toma de
% decisiones en el ámbito de la seguridad ciudadana. Sin embargo, este sistema no se encuentra disponible para el público en general
% por lo cual en \cite{velozSpatialPatternsCrime} se menciona que en ciudades como Quito se han realizado estudios sobre las tendencias
% y patrones delictivos empleando herramientas de análisis espacial mediante sistemas de información geográfica. En la ciudad de Riobamba
% también se han realizado estudios relacionados al análisis de los delitos, en este caso en \cite{escuderoSpatiallyCorrelatedModel2022}
% se propuso estructurar un modelo con el cual realizar un conteo de crímenes basados en la correlación geográfica y temporal de los mismos.

\bigbreak
En la ciudad de Ambato, se evidencia la carencia de sistemas que permitan establecer una fuente confiable de datos georreferenciados para
analizar exhaustivamente los delitos y sus tendencias criminales. El autor de \cite{alvarezCompendioEconomicoSocial2022} determinó, a
través de encuestas a la ciudadanía del cantón Ambato, que el 26.30\% de la población ha sido víctima de actos delictivos, situando a
Ambato en el octavo puesto a nivel cantonal en Ecuador en cuanto al número de denuncias por robos. Según \cite{quintanaAnalisisSistematicoAumento2023},
la cantidad de incidentes registrados en Ambato afecta el nivel de confianza de la ciudadanía, debido al temor que sienten de ser víctimas
de algún incidente y a la preocupación por la integridad de sus familiares.


% Siendo así que más del 92.8\% de la población ambateña se considera propenso a ser víctima de algún delito.

\section{Justificación}
En \cite{benitezalvarezImplementingComprehensiveCitizen2019} se señala que en 2023 Ecuador se
ha visto envuelto en una gran cantidad de incidentes violentos, empezando
desde asaltos, robos, secuestros, violaciones hasta muertes violentas, causando preocupación
en la ciudadanía. En ese contexto, en \cite{alvarezCompendioEconomicoSocial2022} se señala
cómo la mayor cantidad de reportes a la policía proceden de denuncias administrativas.
Estas denuncias no corresponden al total de delitos reales que han ocurrido, debido al hecho de
que no todos los ciudadanos recurren a denunciar formalmente los delitos de los cuales han sido víctimas.
% Tomando en cuenta que desde el año 2008 el Ecuador se ha enfrentado a una
% crisis en relación al crimen e inseguridad llegando a niveles inquietantes de hasta un 18.9\%
% por cada 100.000 habitantes que se han visto afectados por la delincuencia.
Según \cite{castroaniyarPaintingsCrimeComposed2019} esto causa que los datos obtenidos no sean
completamente integrales y precisos en cuanto al lugar y momento en que ocurrieron los incidentes, y
a su vez, dificultan su uso en procesos de análisis y toma de decisiones adecuadas.

\bigbreak
% Según datos registrados por el Compendio económico y social OBEST - INFORMA mediante una encuesta dirigida a todos los ciudadanos del cantón Ambato,
% se determinó que el 26.30\% fue víctima de algún tipo de siniestro entre los períodos de octubre de 2020 y marzo de 2022,
% siendo los delitos más comunes los asaltos y robos con un 33\% del total de delitos y el robo a viviendas con un 8\% del total \cite{alvarezCompendioEconomicoSocial2022}.\\

% En \cite{comisionestadisticadeseguridadciudadanayjusticiaEstadisticasSeguridadIntegral2022} según datos
% obtenidos de la Comisión Estadística de Seguridad Ciudadana y Justicia, entre los años 2021 y 2022,
% se ha notado un incremento de casi un 100\% en el número de homicidios pasando de 2496 a 4823 casos reportados
% y al menos un 25\% en el robo a personas incrementando de 25440 a 31485 incidentes registrados.

% El autor de \cite{naranjo-avalosImpactoGeorreferenciacionColaborativa2019} menciona que se debe
% tomar en cuenta la necesidad de otorgar acceso a la ciudadanía a un sistema que permita la
% recolección de información georreferenciada sobre delitos, dado que la misma información ingresada pueda
% ser validada por la propia ciudadanía, quien, en su defecto, son las víctimas de los siniestros.
% Considerando
% lo expuesto, se justifica la necesidad de crear una herramienta accesible al público y de uso intuitivo. Esta
% plataforma permitirá a los ciudadanos registrar datos georreferenciados sobre los incidentes de los que han sido
% víctimas. Además, brindará un medio para reportar emergencias en tiempo real. Este recurso se convertirá en una
% valiosa fuente de información con la cual se pueda desarrollar un modelo analítico de BI crucial para la toma de
% decisiones que contribuyan significativamente a mejorar la seguridad pública.


El autor de \cite{naranjo-avalosImpactoGeorreferenciacionColaborativa2019} menciona la importancia de proporcionar a
la ciudadanía un sistema accesible que permita recopilar información georreferenciada sobre delitos. Esta información,
una vez ingresada, podría ser validada por los propios ciudadanos, quienes, en su defecto, son las víctimas de estos
incidentes. En este contexto, se justifica la necesidad de desarrollar una herramienta pública, intuitiva y de fácil acceso.
Esta herramienta posibilitaría a los ciudadanos registrar datos georreferenciados  de incidentes que han experimentado
y también ofrecería un canal para reportar emergencias en tiempo real. La creación de esta  herramienta se vuelve crucial,
ya que se convertiría en una valiosa fuente de información. Esta base de datos georreferenciados  podría ser utilizada
para desarrollar un modelo analítico de Business Intelligence (BI) que sería fundamental para la toma de  decisiones
en materia de seguridad pública.

\bigbreak
El presente proyecto es técnicamente factible, ya que las herramientas necesarias para su desarrollo son de acceso libre y
fácil. Además, las especificaciones requeridas del sistema permiten su integración tanto en una aplicación web como móvil.
Desde una perspectiva operativa, el sistema puede ser desarrollado, ya que la recolección de datos será realizada por los
propios usuarios del sistema, quienes son ciudadanos de Ambato.
En cuanto a la factibilidad económica, el proyecto puede llevarse a cabo con costos bajos, los cuales serán asumidos por
el autor de la investigación. Esta estructura financiera asegura la independencia del proyecto al no depender de financiamiento
externo, garantizando su viabilidad económica a lo largo de su desarrollo.

\section{Objetivo general}
Implementar un sistema web y móvil que mejore la recolección de datos
georreferenciados precisos de delitos en la ciudad de Ambato.

% var independiente: sistema web y móvil.
% var dependencia: datos georreferenciados.
% factor medible: cantidad de datos, velocidad de subida, mostrar en tiempo real

\section{Objetivos específicos}
% General:
% Implementar un sistema web y móvil que mejore la recolección de datos
% georreferenciados precisos de delitos con los cuales desarrollar un modelo analítico de
% Business Intelligence (BI) orientado al seguimiento de las tendencias de siniestros
% en la ciudad de Ambato.

\begin{itemize}
      \item Analizar frameworks para el desarrollo web y móvil, enfocados en sistemas para la obtención
            de datos georreferenciados.
      \item Evaluar la precisión de los datos georreferenciados mediante el empleo de técnicas
            estadísticas para determinar la fiabilidad de los datos obtenidos.
            % \item Evaluar la precisión de los datos georreferenciados mediante el empleo de técnicas
            %       de georreferenciación inversa que permitan determinar la exactitud de las coordenadas
            %       obtenidas al verificarlas con la información de su ubicación inversa.
      \item Desarrollar  un sistema web y móvil para la obtención de datos
            georreferenciados de delitos con los cuales desarrollar un modelo analítico de
            Business Intelligence (BI) orientado al seguimiento de las tendencias de siniestros
            en la ciudad de Ambato.
\end{itemize}

\section{Fundamentación teórica}
En este apartado se mostrará los \textit{Antecedentes Investigativos} y \textit{Marco Teórico Conceptual}
resultado de la revisión exhaustiva realizada en los trabajos de autores con investigaciones similares
sobre el tema de investigación.
\subsection{Antecedentes investigativos}
En esta sección se muestran todos aquellos trabajos que proporcionan un contexto
precedente sobre el trabajo de investigación a realizarse.
\bigbreak
Con el propósito de obtener información relevante sobre el tema propuesto para este proyecto se realizó una
búsqueda exhaustiva en las revistas IEEE, SPRINGER LINK, así como en los repositorios de GOOGLE SCHOLAR
y LA UNIVERSIDAD TÉCNICA DE AMBATO en el periodo de 2015 a 2023, en la tabla \ref{table:articulosobtenidos}
se muestran los resultados
obtenidos:
% \begin{ThreePartTable}
%     \begin{TableNotes}[flushleft]
%         \centering
%         \item Elaborado por: el investigador
%     \end{TableNotes}
%     \begin{longtable}{|l|l|l|}
%         \caption{Artículos obtenidos}
%         \label{table:articulosobtenidos}                                                                  \\
%         \hline
%         \textbf{Cadena}                   & \textbf{Revista / Repositorio} & \textbf{Número de artículos} \\
%         \hline
%         (''georeferenced data''           & IEEE                           & 28                           \\ \cline{2-3}
%         OR ''spatiotemporal data''        & SPRINGER LINK                  & 417                          \\ \cline{2-3}
%         OR ''spatio-temporal              & GOOGLE SCHOLAR                 & 1.610                        \\ \cline{2-3}
%         data'' OR ''spatio                & UNIVERSIDAD                    & 0                            \\
%         temporal data'') AND              & TÉCNICA DE AMBATO              &                              \\ \cline{2-3}
%         (''crime analysis''               &                                &                              \\
%         OR crime OR felony OR             &                                &                              \\
%         (''business intelligence''        &                                &                              \\
%         OR BI OR ''analytical             &                                &                              \\
%         ''felony acts'') OR (date         &                                &                              \\
%         OR dates OR time OR  hour         &                                &                              \\
%         OR hours OR ''crime types'')      &                                &                              \\
%         AND  model'' OR ''analytical      &                                &                              \\
%         models'') AND (''crime quantity'' &                                &                              \\
%         OR ''crime count'' OR ''number    &                                &                              \\
%         of crimes'' OR ''total crimes'')  &                                &                              \\
%         \hline
%                                           & Total                          & 2055                         \\
%         \hline
%         \insertTableNotes
%     \end{longtable}
% \end{ThreePartTable}

\begin{longtable}{|l|l|l|}
    \caption{Artículos obtenidos} \label{tab:articulos-obtenidos}                                                                                            \\

    \hline \multicolumn{1}{|c|}{\textbf{Cadena}} & \multicolumn{1}{|c|}{\textbf{Revista / Repositorio}} & \multicolumn{1}{|c|}{\textbf{Número de artículos}} \\ \hline
    \endfirsthead

    \multicolumn{3}{c}%
    {{\normalfont \tablename\ \thetable{} -- continuación de la página anterior}}                                                                            \\
    \hline \multicolumn{1}{|c|}{\textbf{Cadena}} & \multicolumn{1}{|c|}{\textbf{Revista / Repositorio}} & \multicolumn{1}{|c|}{\textbf{Número de artículos}} \\ \hline
    \endhead

    \hline \multicolumn{3}{|r|}{{Continua en la siguiente página}}                                                                                           \\ \hline
    \endfoot

    \hline \hline
    \endlastfoot
    (''georeferenced data''                      & IEEE                                                 & 28                                                 \\ \cline{2-3}
    OR ''spatiotemporal data''                   & SPRINGER LINK                                        & 417                                                \\ \cline{2-3}
    OR ''spatio-temporal                         & GOOGLE SCHOLAR                                       & 1.610                                              \\ \cline{2-3}
    data'' OR ''spatio                           & UNIVERSIDAD                                          & 0                                                  \\ \cline{2-3}
    temporal data'') AND                         & TÉCNICA DE AMBATO                                    &                                                    \\ \cline{2-3}
    (''crime analysis''                          &                                                      &                                                    \\ \cline{2-3}
    OR crime OR felony OR                        &                                                      &                                                    \\ \cline{2-3}
    (''business intelligence''                   &                                                      &                                                    \\ \cline{2-3}
    OR BI OR ''analytical                        &                                                      &                                                    \\ \cline{2-3}
    ''felony acts'') OR (date                    &                                                      &                                                    \\ \cline{2-3}
    OR dates OR time OR  hour                    &                                                      &                                                    \\ \cline{2-3}
    OR hours OR ''crime types'')                 &                                                      &                                                    \\ \cline{2-3}
    AND  model'' OR ''analytical                 &                                                      &                                                    \\ \cline{2-3}
    models'') AND (''crime quantity''            &                                                      &                                                    \\ \cline{2-3}
    OR ''crime count'' OR ''number               &                                                      &                                                    \\ \cline{2-3}
    of crimes'' OR ''total crimes'')             &                                                      &                                                    \\
\end{longtable}
Como se muestra en la tabla \ref{table:articulosobtenidos}, mediante la revisión sistemática se obtuvo un total de 2055 artículos relacionados.
Sin embargo, Con la finalidad de obtener información relevante para el desarrollo de este proyecto se descartaron
aquellos artículos que no cumplan con los criterios de exclusión e inclusión que se muestran en la tabla \ref{table:criteriosinclusionexclusion}:
\begin{ThreePartTable}
    \begin{TableNotes}[flushleft]
        \centering
        \item Elaborado por: el investigador
    \end{TableNotes}
    \begin{longtable}{|l|l|}
        \caption{Criterios de exclusión e Inclusión}
        \label{table:criteriosinclusionexclusion}                                 \\
        \hline
        \textbf{Inclusión}          & \textbf{Exclusión}                          \\ \hline
        \tabitem Crimen             & \tabitem Solo artículos científicos y tesis \\
        \tabitem geointeligencia    & \tabitem Solo textos en español e inglés    \\
        \tabitem geodata            & \tabitem No contiene georreferenciación     \\
        \tabitem BI                 &                                             \\
        \tabitem Analítica          &                                             \\
        \tabitem Georreferenciación &                                             \\
        \hline
        \insertTableNotes
    \end{longtable}
\end{ThreePartTable}

Mediante la exclusión realizada se obtuvo un total de 20 artículos que contiene información sobre datos georreferenciados.
Aunque los artículos obtenidos no cumplen exactamente con el objetivo propuesto para esta investigación, se destaca a continuación
los trabajos más relevantes para el desarrollo de este proyecto.
% Mediante la revision sistemática de la bibliografía en revistas y repositorios digitales, se  no se encontraron trabajos que cumplan exactamente con el objetivo
% propuesto para esta investigación. Sin embargo, se destaca a continuación los trabajos con la información mas relevante para el desarrollo de este proyecto.
\bigbreak
Los autores de \cite{herreraGeoBIBigVGI2015} mencionan las ventajas del uso de los Sistemas de Información Geográfica (GIS).
Estas son herramientas masivas utilizadas para mostrar información de interés en mapas. Además, destacan las ventajas del empleo de
herramientas de Business Intelligence para el procesamiento de grandes cantidades de datos, especialmente en el contexto de
datos georreferenciados, denominado Geo Business Intelligence (GeoBI). Esta disciplina se encarga de abordar el análisis de datos
geoespaciales para mejorar la toma de decisiones en las organizaciones. En el trabajo se propone crear un sistema de Información
Geográfica Voluntaria, del inglés (VGI) el cual conste de un portal web. Este sistema habilitará el manejo de usuarios y visualización de mapas.
También se utilizará una aplicación cliente sea web o móvil, para ingresos de datos y una base de datos geográfica (GeoDB) con
la cual almacenar la información geográfica voluntaria. En el sistema se hará uso del framework "Monitor, Analyze and Drill to Detail"
(MAD), el cual permite separar las necesidades de los usuarios dependiendo al perfil que tengan asignado, como administradores,
analistas y entes de control de una organización. Se menciona que los sistemas GIS deberían proporcionar ciertos elementos básicos,
como portales web, aplicaciones clientes, servidores de mapas y bases de datos con capacidades de almacenamiento de datos geográficos.
Sin embargo, todo lo planteado en el trabajo es una propuesta por parte de los autores y no se muestra un prototipo funcional.
\bigbreak
En \cite{linUsingMachineLearning2017} el trabajo propuesto se centra en el uso de machine learning para la predicción de delitos,
específicamente en el uso de deep learning. En el trabajo se divide un mapa en cuadriculas, con el fin de verificar si en una cuadricula
se a detectado un crimen y tomarlo como un punto caliente. Para este modelo se tomaron en cuenta siete patrones espacio temporales,
las cuales son: mes, mes siguiente, año pasado, cuadrículas rodeadas en ocho direcciones en el mes actual, cuadrículas rodeadas en
ocho direcciones el año pasado, tendencia y proporción, más no se menciona de donde fueron extraídas tales características ni mediante
que medios. No obstante aunque no se trate el tema del software en la adquisición de los datos georreferenciados parad delitos, las
características que se muestran aportan información valiosa para ser utilizada en el contexto de modelos analíticos de BI.
\bigbreak
En \cite{CrimesPredictionUsing2019}, el autor propone un modelo de predicción de delitos basado en el uso de datos espacio-temporales
y el método Kernel Density Estimation (KDE), frecuentemente utilizado para estimar la densidad de crímenes. En el trabajo propuesto,
se detalla el flujo de trabajo para el preprocesamiento de los datos, que implica la eliminación de características irrelevantes y valores
geográficos fuera de la zona de estudio. El autor también menciona que el número de observaciones es de 1,762,311 y detalla las características
del dataset, como fecha, categoría del crimen, día de la semana, distrito policial, resolución del crimen, dirección, latitud y longitud.
A pesar de esto, en el trabajo no se menciona el uso de software para la adquisición de datos georreferenciados, ya que el dataset proviene de Kaggle.
No obstante, el artículo es importante para el desarrollo de este proyecto, ya que ofrece datos prácticos que pueden ser útiles en la creación
de modelos analíticos de Business Intelligence (BI).
\bigbreak
Los autores de \cite{winPCPDParallelCrime2019} proponen el algoritmo Criminal Activity Clustering (CAC) el cual está basado
en la clusterización difusa de puntos delictivos con la finalidad de agrupar grandes volúmenes de datos georreferenciados de delitos,
así como identificar patrones de comportamiento de los siniestros. En el trabajo se muestra una descripción de las características
utilizadas en el desarrollo del algoritmo, las cuales son: fecha, lugar del incidente, información del incidente, información del ataque,
información de la víctima, información del perpetrador, información del arma, causalidades y consecuencias, estas características son
obtenidas desde la base de Datos de Terrorismo Global, del inglés (GTD). Tales características propuestas son útiles para su uso en modelos
analíticos de BI. Sin embargo, el trabajo como tal no aborda el uso de software para la adquisición de datos georreferenciados o
su uso en BI.
\bigbreak
En el trabajo propuesto en \cite{amirkhanyanMethodsFrameworksGeoSpatioTemporal2020} no se aborda el uso de software para la
adquisición de datos georreferenciados de crímenes ni su utilización en modelos de BI. En su lugar, se tratan marcos para el
análisis de datos geoespaciales, los cuales incluyen la recopilación, normalización, geolocalización y almacenamiento de los datos.
Estos marcos permiten crear un algoritmo de clustering del lado del servidor para mapas, capaz de manejar una gran cantidad de datos
georreferenciados de forma masiva. En el trabajo se presentan los diagramas de entidad-relación con los tipos de datos utilizados
para desarrollar el algoritmo.
\bigbreak
En el trabajo presentado en \cite{kumarCrimePredictionUsing2020}, se propone un modelo de predicción de delitos basado en el uso
del algoritmo de K-Nearest Neighboring (KNN). Aunque en el trabajo propuesto no se considera el uso de software para la adquisición de datos
georreferenciados de delitos ni el desarrollo de un modelo analítico de BI, se muestra cómo los datos fueron preprocesados eliminando
valores nulos y datos innecesarios, obteniendo las siguientes características importantes: Hora, latitud, longitud, año, mes y semana del año.
Aun así, el autor menciona que los atributos obtenidos pueden no ser suficientes, por lo cual se deberían encontrar una mayor cantidad
de atributos importantes relacionados con los crímenes.
\bigbreak
En \cite{hossainCrimePredictionUsing2020} el autor se centra principalmente en el uso de algoritmos de aprendizaje supervisado
para la predicción de delitos. La finalidad del uso de estos algoritmos es proporcionar información temprana sobre la actividad
delictiva a entidades de control y agentes de la ley. El trabajo presenta los atributos propuestos para el dataset tales como:
fecha, categoría, día, distrito policial, resolución del crimen, dirección, latitud y longitud. Estos datos son extraídos del
sistema de reportes de incidentes del Departamento de Policía de San Francisco (SFPD). Sin embargo, el trabajo no aborda el uso
de software para la adquisición de datos georreferenciados de crímenes ni su uso para modelos analíticos de Business Intelligence
(BI). A pesar de ello, este trabajo resulta relevante para el desarrollo del proyecto actual, ya que proporciona datos útiles que
pueden ser empleados en la creación de un modelo analítico de BI.
\bigbreak
% El autor en \cite{moralesgutamaSistemaWebUsando2023} presenta el desarrollo de un sistema web el cual permita gestionar las
% calificaciones y asistencia en la unidad educativa Huachi Grande. Este sistema web esta desarrollado tanto su frontend y backend
% utilizando el framework Laravel con el lenguaje de programación PHP, junto a una base de datos en MySQL. Para el desarrollo de
% este proyecto de aplico la metodología de desarrollo ágil XP.
% \bigbreak
En \cite{botto-tobarAppliedTechnologiesSecond2021} se muestra un prototipo de una aplicación web desarrollada
utilizando el patrón de arquitectura cliente-servidor en Visual Studio 2017 Community, utilizando el lenguaje de programación C\#
junto al framework ASP.NET, también la aplicación utiliza el API de google map. Esta aplicación permite recolectar y mostrar datos
georreferenciados de delitos en la ciudad de Ambato. El trabajo propuesto exhibe el modelo de la base de datos, desarrollado en SQL
Server  y las interfaces propuestas para el prototipo. La aplicación desarrollada permite a los usuarios reportar delitos mediante
un formulario y muestra los puntos de siniestros mediante un mapa de calor de la ciudad. Además, permite visualizar reportes de los
delitos desde un punto de vista cualitativo y cuantitativo mediante el uso de la herramienta Power BI integrada en la aplicación web.
No obstante, el trabajo presentado no profundiza en el uso de los datos obtenidos para generar modelos analíticos de BI.
\bigbreak
En \cite{chasichangoAplicacionMovilApoyo2022} se destaca el uso de la georreferenciación en sistemas de seguridad en el cual
se propone una aplicación móvil desarrollada en IONIC usando el lenguaje de programación Javascript. La aplicación permite el
envío de geolocalización mediante notificaciones push utilizando el servicio de one signal. El backend de esta aplicación esta
desarrollado en NodeJS utilizando express y como base de datos utiliza MongoDB. El objetivo de esta aplicación es ayudar acometer
la problemática social de la inseguridad ciudadana en la parroquia Santa Rosa de la ciudad de Ambato.
\bigbreak
En \cite{lesanoperezAplicativoMovilGeoubicacion2022} se describe el desarrollo de una aplicación móvil que se encarga de enviar
la ubicación en tiempo real para el transporte privado de la empresa Plasticaucho S.A. de la ciudad de Ambato. Esta aplicación móvil
utiliza Xamarin Forms con el lenguaje de programación C\#. En el backend, se empleó el framework ASP.NET, también utilizando el lenguaje
de programación C\#, e implementando el patrón de arquitectura Model View Controller (MVC), junto a una base de datos en PostgreSQL. El proyecto se llevó a
cabo utilizando la metodología de desarrollo ágil Extreme programming (XP).
\bigbreak
En \cite{chicaizavillegasAplicacionWebPara2023}, se desarrolló una aplicación web que incorpora georreferenciación.
Esta aplicación utiliza el framework Angular junto con el lenguaje de programación JavaScript para el frontend. El backend fue
implementado mediante el framework ASP.NET utilizando el lenguaje de programación C\#, y se empleó una base de datos MySQL. Para
el desarrollo de este proyecto, se aplicó la metodología de desarrollo Extreme Programming (XP), orientada a grupos reducidos de
desarrollo. La aplicación creada permite gestionar servicios, notificaciones de pago y geolocalización de viviendas para la empresa
Optynet.
\bigbreak
Los autores en \cite{cantilloAplicativoMovilPara} presentan el diseño y desarrollo de una aplicación móvil que aborda la
problemática de la predicción de crímenes. Esta aplicación se fundamenta en la integración de tecnologías modernas, la recopilación
de datos geoespaciales, el análisis de datos e inteligencia artificial. En el trabajo, se detallan las técnicas utilizadas para la
recopilación de información. Se menciona el empleo de la herramienta Taled Open Studio, la cual se utilizó para llevar a cabo el
proceso de Extracción, Transformación y Carga de datos (ETL), almacenándolos en una base de datos MySql. Además, el proyecto exhibe
las interfaces de la aplicación móvil, que constan de un módulo de autenticación, otro de reporte de delitos y un tercer módulo para
la visualización de estos. Dentro del módulo de visualización de delitos, se presentan mapas de calor que muestran los incidentes
ocurridos, agrupándolos por áreas mediante esferas de colores que representan el total de siniestros en cada zona. Por otro lado,
el módulo de reporte de delitos permite a los usuarios informar los incidentes de los que han sido víctimas mediante un formulario
y un mapa para señalar la ubicación del suceso. Además, la aplicación ofrece estadísticas sobre los incidentes reportados en distintos
periodos de tiempo para el conocimiento de la ciudadanía. Este trabajo proporciona un medio para la adquisición de datos georreferenciados
mediante el uso de software, así como una plataforma para mostrar información relevante a la ciudadanía. No obstante, no incluye detalles
sobre las características utilizadas ni los diagramas de la base de datos. Tampoco aborda el uso de modelos analíticos de Business
Intelligence (BI).

\subsection{Marco teórico}
Esta sección establecerá las bases, fundamentos y teorías pertinentes a través de una exhaustiva revisión de
antecedentes. Esto permitirá respaldar y contextualizar el desarrollo de la presente investigación, ofreciendo
un marco sólido que sustente cada paso del estudio.

\subsubsection{Metodología ágil}
Una metodología ágil es un conjunto de técnicas aplicadas utilizadas con el propósito de facilitar los tiempos de
desarrollo y obtener un producto final de calidad. Los autores de \cite{lesanoperezAplicativoMovilGeoubicacion2022} y
\cite{chicaizavillegasAplicacionWebPara2023}
señalan el uso de la
metodología ágil Extreme Programming (XP). Los autores
mencionan la flexibilidad y facilidad de modelado que ofrece la metodología ágil XP, además, de mejorar los tiempos
de desarrollo, calidad del producto y ayuda a minimizar los errores durante el desarrollo. También
se indica que la metodología XP está orienta a grupos de
trabajo pequeño, además, de adaptarse mejor a las necesidades de los proyectos ya que es versátil al incorporar
cambios durante el desarrollo. Así mismo los autores señalan que una de las
principales ventajas de la metodología XP es la capacidad de dividir el trabajo en pequeñas historias de usuario
las cuales son completamente funcionales, lo cual garantiza la entrega
del producto final en etapas manejables.

\bigbreak
En \cite{chasichangoAplicacionMovilApoyo2022} se utilizó la metodología de desarrollo Mobile-D, ya que combina
beneficios de metodologías como XP, Crystal y RUP. El autor menciona que la metodología Mobile-D se enfoca en el
desarrollo de aplicaciones móviles, especialmente útil en aplicaciones pequeñas con pocos módulos y vistas. Además,
ofrece un estilo de desarrollo por iteraciones, el cual es rápido y útil para grupos pequeños.

\subsubsection{Framework Web}
Un framework web constituye un conjunto de herramientas y librerías esenciales para el desarrollo de aplicaciones web.
En \cite{chicaizavillegasAplicacionWebPara2023}, el autor destaca Angular como un framework que posibilita el diseño
de sistemas web estables y escalables. Esto se debe a las capacidades que ofrece para crear aplicaciones SPA eficientes.
Asimismo, menciona que Angular proporciona una arquitectura basada en componentes y módulos organizados, los cuales se
integran fácilmente con otras tecnologías, como el caso de ASP.NET Core, utilizado como plantilla de interfaz de usuario.
% Por otro lado, en \cite{moralesgutamaSistemaWebUsando2023}, el autor describe a Laravel como un framework web que ofrece
% una arquitectura limpia y cuenta con una amplia comunidad de desarrolladores. Esto asegura un soporte extenso y una fuerte
% presencia en el mercado. Además, señala que Laravel dispone de una gran cantidad de herramientas y bibliotecas que simplifican
% la implementación de diversas funcionalidades.

\subsubsection{Framework Backend}
Un framework backend consiste en un conjunto de herramientas y librerías predefinidas que permiten desarrollar y mantener
la lógica de una aplicación en el lado del servidor. En \cite{chicaizavillegasAplicacionWebPara2023}, el autor menciona
que ASP.NET es un framework backend gratuito, seguro, de alto rendimiento y de código abierto. El autor destaca que ASP.NET
permite ejecutar aplicaciones web en sistemas operativos como Windows, macOS y Linux utilizando C\# en el lado del servidor,
% además de otras tecnologías en el lado del cliente. Por otra parte, en \cite{moralesgutamaSistemaWebUsando2023}, el autor indica que Laravel,
% como framework backend, es robusto y completo. Se mencionan también las medidas de seguridad incluidas en Laravel para
% proteger las aplicaciones contra ataques de inyección SQL y XSS. El autor explica que Laravel cuenta con un Object-Relational
% Mapping (ORM) llamado Eloquent, que permite manipular la base de datos de manera intuitiva y eficiente. En cambio, 
el autor de
\cite{chasichangoAplicacionMovilApoyo2022} destaca a Node.js como tecnología backend  la cual permitiendo el manejo eficiente
de múltiples procesos simultáneamente. Además, se menciona que la curva de aprendizaje de Node.js es baja, al igual que su
costo de desarrollo. También se resalta que Node.js ofrece una manera sencilla de crear APIs REST utilizando Express,
las API REST permiten comunicar diferentes sistemas mediante los protocolos HTTP, lo que facilita la creación rápida y
flexible de un backend completo.

\subsubsection{Aplicación Movil}
Una aplicación móvil se refiere a un tipo de aplicación diseñada para ejecutarse en teléfonos inteligentes o tabletas.
En \cite{chasichangoAplicacionMovilApoyo2022}, el autor destaca que IONIC, utilizado en el desarrollo de aplicaciones
móviles, presenta una curva de aprendizaje reducida, ya que puede combinarse con frameworks como Angular y librerías de
componentes como React y Svelte. Además, menciona que IONIC permite la creación de componentes reutilizables, lo que conlleva
a un código más legible y fácil de mantener. Asimismo, resalta su extensa documentación respaldada por una gran comunidad,
lo que facilita la solución a errores de manera más sencilla. Por otra parte, en \cite{lesanoperezAplicativoMovilGeoubicacion2022},
el autor enfatiza el uso de Xamarin Forms para el diseño de aplicaciones móviles debido a su capacidad para crear aplicaciones
multiplataforma. Expone que Xamarin permite el desarrollo de aplicaciones utilizando el lenguaje C\#, el cual se compila en
código nativo para Android e iOS. Además, al ser una tecnología de Microsoft, posibilita la integración sencilla de paquetes
externos mediante NuGet.

\subsubsection{Bases de datos}
Una base de datos es un sistema electrónico que permite el almacenamiento de datos de forma organizada.
En \cite{chicaizavillegasAplicacionWebPara2023}, el autor destaca que MySQL, como Sistema Gestor de Bases de Datos Relacionales
(SGBDR), ofrece un alto rendimiento y fiabilidad, lo que permite almacenar, buscar, ordenar y obtener información de forma eficiente.
% Por otro lado, en \cite{moralesgutamaSistemaWebUsando2023}, 
También se indica que MySQL es uno de los SGBDR más populares, siendo de
código abierto y respaldado por Oracle. En el mismo sentido, en \cite{lesanoperezAplicativoMovilGeoubicacion2022}, el autor
resalta a PostgreSQL como SGBDR por ser uno de los gestores más rápidos y seguros del mundo. Expone su gran escalabilidad,
facilidad de uso y capacidad para manejar no solo variables primitivas, sino también variables tipo objeto. Además, señala que
la base de datos de PostgreSQL es ilimitada, permitiendo un número ilimitado de campos por tabla.

\bigbreak
En otro enfoque, el autor de \cite{chasichangoAplicacionMovilApoyo2022} menciona a MongoDB como una base de datos NoSQL capaz
de manejar una gran cantidad de información. Se destaca su carácter de código abierto y su flexibilidad para trabajar con
diversas tecnologías de backend. Además, se resalta su escalabilidad al estar vinculada a servidores de Amazon, Google y
Microsoft, ofreciendo también planes gratuitos convenientes para aplicaciones pequeñas.

\subsubsection{Mapa web}
Un mapa web es una visualización interactiva de información geográfica que se puede abrir en navegadores y dispositivos móviles.
En \cite{botto-tobarAppliedTechnologiesSecond2021}, el autor destaca que el API proporcionado por Google Maps para renderizar
mapas ofrece varias características y prestaciones, siendo uno de los productos más importantes de Google. También menciona que
una de las principales ventajas de esta API es su capacidad para integrarse con distintos lenguajes de programación y, por ende,
con diferentes frameworks. Esto hace que sea altamente flexible para su uso en diferentes sistemas. Por otro lado, en
\cite{gomezcantilloAplicativoMovilPara}, el autor señala que Leaflet permite crear mapas interactivos y aplicaciones de
mapas en línea, además de ser de código abierto. También destaca la amplia gama de herramientas disponibles en Leaflet que
facilitan la visualización de datos georreferenciados. Incluso, ofrece a los desarrolladores la capacidad de crear experiencias
de usuario altamente interactivas y personalizables, lo cual resulta ideal para aplicaciones tanto móviles como web.

\subsubsection{Geo Business Intelligence (GEOBI)}
En \cite{herreraGeoBIBigVGI2015} se empleó GEOBI dado que esta mejora la toma de decisiones en organizaciones. El
GEOBI proporciona información más precisa y oportuna incorporando la dimensión espacial en la información. También
permite que sus componentes se vuelvan geográficos, es decir, en lugar de OLAP se utiliza SOLAP una combinación de un
Sistema de información geográfica (SIG) y un cubo OLAP, además, sus bases de datos se convierten en GeoDB, bases de
datos que contienen información geográfica.

\subsubsection{Patron MVC}
En \cite{lesanoperezAplicativoMovilGeoubicacion2022}, el autor menciona la importancia del patrón Model View Controller
(MVC), el cual permite separar una aplicación en tres capas de componentes: Modelo, Vista y Controlador. El autor señala que
la separación de funciones de este patrón de arquitectura mejora la mantenibilidad del código y hace más sencillo depurar
errores. El autor indica que el uso del patrón MVC permite disminuir los errores, dado que las dependencias no están
distribuidas en diferentes secciones del código y, por lo tanto, evita mezclar la lógica de negocio con las interfaces
de usuario.

\subsubsection{Diagrama UML}
En los artículos revisados no se señala el uso de diagramas UML en el desarrollo de sistemas de software. Sin embargo, se
considera necesario mencionarlo para el desarrollo de este proyecto. El lenguaje de modelado unificado (UML) es una
representación visual y semántica de la arquitectura, diseño e implementación de sistemas de software. Los diagramas
UML permiten describir los límites, la estructura, el comportamiento y los objetos que contiene un sistema. Aunque este
no es un lenguaje de programación, mediante herramientas externas puede ser utilizado para la generación de código para
diferentes lenguajes. Según el Object Management Group (OMG), el propósito de UML es proveer a los ingenieros,
desarrolladores y arquitectos de software un medio por el cual puedan realizar análisis, diseño e implementación de
sistemas de software. Aunque UML no está enfocado en bases de datos, este permite crear modelos de datos conceptuales
de alto nivel.

\subsubsection{NextJS}
En los artículos revisados no se menciona el uso de NextJS como framework frontend para el desarrollo de aplicaciones web.
Sin embargo, se considera necesario mencionarlo para el desarrollo de este proyecto. NextJS es un framework web el cual
utiliza la librería de componentes React la cual es desarrollada por Meta. NextJS permite crear aplicaciones web
completas utilizando React para crear interfaces de usuario dinámicas y rápidas basadas en componentes y añadiendo
características adicionales y optimizaciones tales como enrutador basado en archivos, compilación y empaquetado de la
aplicación, optimización de images, Renderizado del lado del servidor (SSR), Renderizado del lado del cliente (CSR),
Generación de sitios estáticos (SSG). NextJS permite desarrollar un backend basado en NodeJS con el cual crear un API
REST en la misma aplicación.

\subsubsection{NestJS}
En los artículos revisados no se menciona el uso de NestJS como framework backend para el desarrollo de API REST.
Sin embargo, se considera necesario mencionarlo para el desarrollo de este proyecto. NestJS es un framework backend
basado en NodeJS el cual permite diseñar aplicaciones en el lado del servidor como. NestJS combina elementos de la
programación orientada a objetos (POO), programación funcional (PF) y programación reactiva funcional junto con con
TypeScript para diseñar aplicaciones eficientes y escalables. NestJS brinda un nivel de abstracción a framework de
servidor HTTP sólidos como Express y Fastify, además de exponer su API directamente con lo cual permite a los
desarrolladores implementar modules externos disponibles para la plataforma.

\subsubsection{Flutter}
En los artículos revisados, no se menciona el uso de Flutter como framework para el desarrollo de aplicaciones móviles;
sin embargo, se considera necesario mencionarlo para el desarrollo de este proyecto. Flutter es un framework desarrollado
por Google para el desarrollo de aplicaciones móviles, web y de escritorio de forma nativa con solo código base. Permite
crear interfaces de usuario visualmente atractivas al manejar los elementos de la interfaz en componentes más pequeños
llamados widgets. Flutter tiene un alto rendimiento, dado que utiliza Dart como lenguaje de programación, el cual compila
código a máquina, permitiendo que las aplicaciones desarrolladas sean rápidas y eficaces. Además, Flutter ofrece una buena
experiencia de desarrollo al incluir herramientas que permiten realizar actualizaciones en caliente y revisar el estado de
los componentes de la aplicación.

\section{Metodología}
\bigbreak

En el proyecto propuesto para las aplicaciones cliente, se empleará una aplicación móvil, dado que esta puede ser utilizada en una gran variedad de
dispositivos, tanto Android como iOS. La aplicación móvil empleará una arquitectura Model View Controller (MVC), con la cual se separará la lógica de
negocio de la interfaz de usuario y la interacción. También se utilizará un framework web, dado que este ofrece una estructura de trabajo establecida,
además de incluir características extras y optimizaciones que facilitan el desarrollo de aplicaciones web. El sistema web se desarrollará utilizando el
patrón de arquitectura MVC, dado que este permite dividir la lógica y responsabilidad en 3 capas y así mantener una mejor estructura en el código.
Para el desarrollo del backend, se emplearán frameworks que agilicen la creación de un API REST eficiente y escalable, el cual permita la comunicación
entre diferentes sistemas, como aplicaciones web y móviles, mediante el protocolo HTTP. Para el almacenamiento de información, se empleará una base de
datos relacional, dado que esta ofrece un gran rendimiento y fiabilidad al manejar grandes volúmenes de datos. Además, permite manejar relaciones entre
diferentes tablas, las cuales ayudan a recuperar datos de forma eficiente, garantizando la integridad referencial. Para el desarrollo del proyecto, se
utilizará una metodología de desarrollo ágil por la facilidad, flexibilidad y versatilidad que ofrece al realizar cambios durante el desarrollo. Además,
mejora los tiempos de desarrollo y la calidad del producto final.

\subsection{Modalidad de investigación}
En esta sección se muestran los enfoques a los que se encuentra orientada la investigación propuesta para este proyecto. A continuación se detalla
cada una de ellas.
\bigbreak
\textbf{Investigación cuantitativa}\\

En el presente proyecto se emplea una metodología de investigación cuantitativa. Esta se centra en la recopilación de
datos numéricos y estadísticos que permiten establecer las tendencias de siniestralidad a través del análisis del número
de crímenes. Esta información se utiliza como base fundamental para el diseño de un modelo analítico de BI.

\bigbreak
\textbf{Investigación aplicada}\\
Se emplea una investigación aplicada, ya que el desarrollo del proyecto propuesto utilizará la teoría, métodos,
técnicas, tecnologías y conocimientos adquiridos a lo largo de ciclos académicos previos. Este enfoque permitirá
aplicar de manera práctica y específica el conocimiento acumulado durante la trayectoria educativa.

% \bigbreak
% \textbf{Investigación bibliográfica}\\
% La Investigación es Bibliográfica y Documental porque es necesario recopilar información de
% documentos como artículos académicos, tesis y libros que sirvan de apoyo
% para la contextualización de la propuesta a desarrollar.

% \bigbreak
% \textbf{Investigación de campo}\\
% La Investigación es de Campo ya que la información y características del problema
% serán extraídas mediante el contacto directo en el lugar de los hechos, es decir, el
% Centro de Transferencia y Desarrollo de Tecnologías y la Dirección de Vinculación
% con la Sociedad.

% \bigbreak
% \textbf{Investigación aplicada}\\
% La Investigación es Aplicada debido a que se emplearán conocimientos
% adquiridos a lo largo de la carrera para el desarrollo de la propuesta.

\subsection{Población y muestra}

Al no contar con una estimación adecuada sobre la cantidad de personas requerida para estipular
la población, se optó por definir una muestra infinita con población desconocida.
\begin{itemize}
    \item n = tamaño de la muestra
    \item Z = nivel de confianza, un nivel de confianza habitual es de 95\% con una puntuación estándar de 1.96
    \item p = probabilidad de éxito o proporción esperada
    \item q = probabilidad de fracaso
    \item pq = varianza de la población, cada una equivale a 0.50, (p=q=0.50)
    \item e = error de estimación máximo aceptable, el error habitual es de 5\% con una puntuación estándar de 0.05
\end{itemize}
\[
    n=\frac{Z^2 \cdot p \cdot q}{e^2}
\]\\
\[
    n=\frac{(1.96)^2 \cdot 0.50 \cdot 0.50}{(0.05)^2}
\]\\
\[
    n=384.16 \approx 385
\]
Con un nivel de confianza del 95\% y un error de estimación máximo aceptable del 5\%, se obtiene una muestra aproximada de 385 personas. \\ \\
\subsection{Recolección de información}
La recolección de información para el presente proyecto se llevará a cabo a través de un
formulario integrado en una aplicación web y móvil. Dicho formulario recopilará los detalles
de las víctimas de los incidentes en la ciudad de Ambato, permitiendo el registro de
información relevante, como el tipo de delito, la fecha y hora del suceso, además de
emplear un mapa para obtener la ubicación geográfica precisa del lugar donde ocurrió el
incidente. También se obtendrán datos por medio de alertas de peligro que podrán ser enviadas
mediante la aplicación móvil.
\subsection{Procesamiento y análisis de datos}
Utilizando los datos recolectados a través del formulario y las alertas de peligro, se
emplearán algoritmos de aumento de datos. El objetivo es obtener una cantidad considerable
de elementos con los que trabajar. Esta información será fundamental para la creación de un
cubo de Procesamiento Analítico en Línea (OLAP), que posibilitará el análisis exhaustivo de
dichos datos y, a su vez, permitirá la generación de informes detallados.
\input{sections/6-metodologia/propuesta-de-solución.tex}
% \subsection{Desarrollo del proyecto}
\begin{itemize}
      \item Analizar frameworks para el
            desarrollo web y móvil, enfocados
            en sistemas para la obtención de
            datos georreferenciados.
            \begin{itemize}
                  \item Elaboración del escrito de la tesis.
                  \item Investigación de tecnologías de
                        desarrollo web.
                  \item Investigación de tecnologías de
                        desarrollo móvil.
                  \item Investigación de tecnologías
                        para mapas.
                  \item Investigación de tecnologías
                        para backend.
                  \item Investigación de metodologías
                        de desarrollo.
                  \item Investigación de gestores de
                        bases de datos.
            \end{itemize}
      \item Evaluar la precisión de los datos
            georreferenciados mediante el
            empleo de técnicas estadísticas
            para determinar la fiabilidad de los
            datos obtenidos.
            \begin{itemize}
                  \item Recopilación de datos
                        georreferenciados.
                  \item Aplicación de técnicas
                        estadísticas para la evaluación
                        de precisión.
                  \item Análisis de resultados y
                        determinación de fiabilidad
                        de los datos.
            \end{itemize}
      \item Desarrollar un sistema web y
            móvil para la obtención de datos
            georreferenciados de delitos con
            los cuales desarrollar un modelo
            analítico de Business Intelligence
            (BI) orientado al seguimiento de
            las tendencias de siniestros en la
            ciudad de Ambato.
            \begin{itemize}
                  \item Diseño del sistema web.
                  \item Diseño del sistema móvil.
                  \item Desarrollo backend.
                  \item Desarrollo del sistema web.
                  \item Desarrollo del sistema móvil.
                  \item Implementación del modelo de
                        Business Intelligence.
                  \item Pruebas y validación del sistema.
                  \item Presentación del trabajo final.
            \end{itemize}

\end{itemize}

\section{Recursos}
En esta sección se especificarán los elementos requeridos para el desarrollo del presente proyecto.
Se muestran aspectos importantes como el equipo necesario, materiales, tecnologías y cualquier otro
elemento vital para llevar a cabo las actividades programadas.
\subsection{Institucionales}
Para el desarrollo de esta investigación, no se cuentan con recursos institucionales, ya que el proyecto
carece de relación con cualquier institución o empresa. Esto implica la ausencia de respaldo financiero,
infraestructura o apoyo logístico proveniente de entidades institucionales o empresariales.

\subsection{Humanos}
\begin{itemize}
    \item Usuarios de la aplicación
\end{itemize}

\subsection{Materiales}
\begin{itemize}
    \item Laptop.
    \item Resma de papel bond.
    \item Impresora.
    \item Tinta.
    \item Material de escritorio.
    \item Smartphone.
    \item Internet.
    \item Servicio de hosting web.
    \item Servicio de hosting backend.
    \item Servicio de hosting BD.
    \item Herramientas de desarrollo (IDE's).
\end{itemize}

\newpage
\subsection{Económicos}
\begin{ThreePartTable}
	\begin{TableNotes}[flushleft]
		\centering
		\item Elaborado por: el investigador
	\end{TableNotes}
	\begin{longtable}{|l |p{0.3\textwidth}| l| l| p{0.12\textwidth}| p{0.12\textwidth}|}
		\caption{Presupuesto}
		\label{table:presupuesto}                                                                                                                                                     \\
		\hline
		\textbf{N.}                                       & \textbf{Detalle}                   & \textbf{Unidad} & \textbf{Cantidad} & \textbf{Valor Unitario} & \textbf{Valor Total} \\
		\hline
		1                                                 & Resma de papel bond                & c/u             & 1                 & \$05.00                 & \$05.00              \\\hline
		2                                                 & Impresora                          & c/u             & 1                 & \$200.00                & \$200.00             \\\hline
		2                                                 & Material de escritorio             & c/u             & 1                 & \$05.00                 & \$05.00              \\\hline
		2                                                 & Herramientas de desarrollo (IDE's) & c/u             & 1                 & \$00.00                 & \$00.00              \\\hline
		2                                                 & Tinta                              & c/u             & 1                 & \$10.00                 & \$10.00              \\\hline
		3                                                 & Laptop                             & c/u             & 1                 & \$950.00                & \$950.00             \\\hline
		4                                                 & Smartphone                         & c/u             & 1                 & \$350.00                & \$350.00             \\\hline
		5                                                 & Servicio de internet               & c/u             & 4                 & \$25.00                 & \$100.00             \\\hline
		6                                                 & Servicio de hosting web            & c/u             & 1                 & \$20.00                 & \$20.00              \\\hline
		7                                                 & Servicio de hosting api            & c/u             & 1                 & \$20.00                 & \$20.00              \\\hline
		8                                                 & Servicio de hosting BD             & c/u             & 1                 & \$20.00                 & \$20.00              \\\hline
		8                                                 & Personas (Usuarios)                & c/u             & 1                 & \$00.00                 & \$00.00              \\
		\hline
		\multicolumn{5}{|l|}{\textbf{Subtotal}}           & \$1680                                                                                                                    \\
		\multicolumn{5}{|l|}{\textbf{Imprevistos (\%15)}} & \$252                                                                                                                     \\
		\multicolumn{5}{|l|}{\textbf{Total}}              & \$1932                                                                                                                    \\
		\hline
		\insertTableNotes
	\end{longtable}
\end{ThreePartTable}


\begin{landscape}
    \section{Cronograma}
    En esta sección se detalla el plan de actividades que se llevarán a cabo durante la ejecución del proyecto de
    investigación. El presente proyecto se realizara desde el 11 de marzo de 2024 hasta el 28 de junio de 2024,
    abarcando desde el inicio hasta la culminación de las tareas y etapas planificadas.
    \begin{ThreePartTable}
        \begin{TableNotes}[flushleft]
            \centering
            \item Elaborado por: el investigador
        \end{TableNotes}
        \begin{longtable}{|p{0.41\textwidth}|p{5.5cm}|p{4cm}|l|l|l|l|l|l|l|l|l|l|l|l|l|l|l|l|l|l|l|}
            \caption{Cronograma}
            \label{table:cronograma}                                                                                                                                                                                                                                                                                                                                                                                                                                                                                                                                                                                                               \\
            \hline
            \multirow{2}{*}{\bfseries Objetivos}                                & \multirow{2}{*}{\bfseries Actividades}                              & \multirow{2}{*}{\bfseries Recursos}                            & \multicolumn{3}{|c|}{\bfseries MARZO} & \multicolumn{4}{|c|}{\bfseries ABRIL} & \multicolumn{4}{|c|}{\bfseries MAYO} & \multicolumn{4}{|c|}{\bfseries JUNIO}                                                                                                                                                                                                                                                              \\
            \cline{4-21}
                                                                                &                                                                     &                                                                & 2                                     & 3                                     & 4                                    & 1                                     & 2                    & 3                    & 4                    & 1                    & 2                    & 3                    & 4                    & 1                    & 2                    & 3                    & 4                    \\
            \hline
            \multirow{6}{0.41\textwidth}{Analizar frameworks para el desarrollo web y móvil, enfocados en sistemas para la obtención
            de datos georreferenciados.}                                        & Elaboración del escrito de la tesis                                 &                                                                & \cellcolor{green!25}                  & \cellcolor{green!25}                  & \cellcolor{green!25}                 & \cellcolor{green!25}                  & \cellcolor{green!25} & \cellcolor{green!25} & \cellcolor{green!25} & \cellcolor{green!25} & \cellcolor{green!25} & \cellcolor{green!25} & \cellcolor{green!25} & \cellcolor{green!25} & \cellcolor{green!25} & \cellcolor{green!25} &                      \\
            \cline{2-2}\cline{4-21}
                                                                                & Investigación de tecnologías de desarrollo web                      & Repositorios digitales                                         & \cellcolor{green!25}                  &                                       &                                      &                                       &                      &                      &                      &                      &                      &                      &                      &                      &                      &                      &                      \\
            \cline{2-2}\cline{4-21}
                                                                                & Investigación de tecnologías de desarrollo móvil                    &                                                                & \cellcolor{green!25}                  &                                       &                                      &                                       &                      &                      &                      &                      &                      &                      &                      &                      &                      &                      &                      \\
            \cline{2-2}\cline{4-21}
                                                                                & Investigación de tecnologías para mapas                             &                                                                &                                       & \cellcolor{green!25}                  &                                      &                                       &                      &                      &                      &                      &                      &                      &                      &                      &                      &                      &                      \\
            \cline{2-2}\cline{4-21}
                                                                                & Investigación de tecnologías para backend                           &                                                                &                                       & \cellcolor{green!25}                  &                                      &                                       &                      &                      &                      &                      &                      &                      &                      &                      &                      &                      &                      \\
            \cline{2-2}\cline{4-21}
                                                                                & Investigación de metodologías de desarrollo                         &                                                                &                                       & \cellcolor{green!25}                  &                                      &                                       &                      &                      &                      &                      &                      &                      &                      &                      &                      &                      &                      \\

            \cline{2-2}\cline{4-21}
                                                                                & Investigación de gestores de bases de datos                         &                                                                &                                       & \cellcolor{green!25}                  &                                      &                                       &                      &                      &                      &                      &                      &                      &                      &                      &                      &                      &                      \\
            \hline
            \multirow{4}{0.41\textwidth}{Evaluar la precisión de los datos georreferenciados mediante el empleo de técnicas
            estadísticas para determinar la fiabilidad de los datos obtenidos.} & Recopilación de datos georreferenciados                             & Formulario para registrar delitos mediante georreferenciación. &                                       &                                       & \cellcolor{green!25}                 & \cellcolor{green!25}                  & \cellcolor{green!25} & \cellcolor{green!25} & \cellcolor{green!25} & \cellcolor{green!25} & \cellcolor{green!25} & \cellcolor{green!25} & \cellcolor{green!25} &                      &                      &                      &                      \\
            \cline{2-2}\cline{4-21}
                                                                                & Aplicación de técnicas estadísticas para la evaluación de precisión & y Técnicas estadísticas                                        &                                       &                                       &                                      &                                       &                      &                      &                      &                      &                      &                      &                      & \cellcolor{green!25} &                      &                      &                      \\
            \cline{2-2}\cline{4-21}
                                                                                & Análisis de resultados y determinación de fiabilidad de los datos   &                                                                &                                       &                                       &                                      &                                       &                      &                      &                      &                      &                      &                      &                      & \cellcolor{green!25} &                      &                      &                      \\
            \hline
            \multirow{4}{0.41\textwidth}{Desarrollar  un sistema web y móvil para la obtención de datos
                georreferenciados de delitos con los cuales desarrollar un modelo analítico de
                Business Intelligence (BI) orientado al seguimiento de las tendencias de siniestros
            en la ciudad de Ambato.}                                            & Diseño del sistema web                                              & Mockups,                                                       &                                       & \                                     & \cellcolor{green!25}                 & \cellcolor{green!25}                  &                      &                      &                      &                      &                      &                      &                      &                      &                      &                      &                      \\
            \cline{2-2}\cline{4-21}
                                                                                & Diseño del sistema móvil                                            & Diagramas UML                                                  &                                       &                                       & \cellcolor{green!25}                 & \cellcolor{green!25}                  &                      &                      &                      &                      &                      &                      &                      &                      &                      &                      &                      \\
            \cline{2-2}\cline{4-21}
                                                                                & Desarrollo backend                                                  & Framework backend,                                             &                                       &                                       &                                      &                                       & \cellcolor{green!25} & \cellcolor{green!25} & \cellcolor{green!25} & \cellcolor{green!25} & \cellcolor{green!25} & \cellcolor{green!25} & \cellcolor{green!25} &                      &                      &                      &                      \\
            \cline{2-2}\cline{4-21}
                                                                                & Desarrollo del sistema web                                          & Framework web,                                                 &                                       &                                       &                                      &                                       & \cellcolor{green!25} & \cellcolor{green!25} & \cellcolor{green!25} & \cellcolor{green!25} &                      &                      &                      &                      &                      &                      &                      \\
            \cline{2-2}\cline{4-21}
                                                                                & Desarrollo del sistema móvil                                        & Framework móvil y                                              &                                       &                                       &                                      &                                       &                      &                      &                      &                      & \cellcolor{green!25} & \cellcolor{green!25} & \cellcolor{green!25} &                      &                      &                      &                      \\
            \cline{2-2}\cline{4-21}
                                                                                & Implementación del modelo de Business Intelligence                  & GEO Business Intelligence                                      &                                       &                                       &                                      &                                       &                      &                      &                      &                      &                      &                      &                      &                      & \cellcolor{green!25} & \cellcolor{green!25} &                      \\
            \cline{2-2}\cline{4-21}
                                                                                & Pruebas y validación del sistema                                    &                                                                &                                       &                                       &                                      &                                       &                      &                      &                      &                      &                      &                      &                      &                      & \cellcolor{green!25} & \cellcolor{green!25} &                      \\
            \cline{2-2}\cline{4-21}
                                                                                & Presentación del trabajo final                                      &                                                                &                                       &                                       &                                      &                                       &                      &                      &                      &                      &                      &                      &                      &                      &                      &                      & \cellcolor{green!25} \\
            \hline
            \insertTableNotes
        \end{longtable}
    \end{ThreePartTable}
\end{landscape}

% \begin{landscape}
%     \begin{ThreePartTable}
%         \begin{TableNotes}[flushleft]
%             \centering
%             \item Elaborado por: el investigador
%         \end{TableNotes}
%         \begin{longtable}{|p{0.41\textwidth}|p{8cm}|l|l|l|l|l|l|l|l|l|l|l|l|l|l|l|l|l|l|l|}
%             \caption{Cronograma}
%             \label{table:cronograma}                                                                                                                                                                                                                                                                                                                                                                                                                                                                                                                                              \\
%             \hline
%             \multirow{2}{*}{\bfseries Objetivos}                                & \multirow{2}{*}{\bfseries Actividades}                              & \multicolumn{3}{|c|}{\bfseries MARZO} & \multicolumn{4}{|c|}{\bfseries ABRIL} & \multicolumn{4}{|c|}{\bfseries MAYO} & \multicolumn{4}{|c|}{\bfseries JUNIO}                                                                                                                                                                                                                                                              \\
%             \cline{3-21}
%                                                                                 &                                                                     & 2                                     & 3                                     & 4                                    & 1                                     & 2                    & 3                    & 4                    & 1                    & 2                    & 3                    & 4                    & 1                    & 2                    & 3                    & 4                    \\
%             \hline
%             \multirow{4}{0.41\textwidth}{Evaluar la precisión de los datos georreferenciados mediante el empleo de técnicas
%             estadísticas para determinar la fiabilidad de los datos obtenidos.} & Recopilación de datos georreferenciados                             &                                       & \cellcolor{green!25}                  & \cellcolor{green!25}                 & \cellcolor{green!25}                  & \cellcolor{green!25} & \cellcolor{green!25} & \cellcolor{green!25} & \cellcolor{green!25} & \cellcolor{green!25} & \cellcolor{green!25} & \cellcolor{green!25} & \cellcolor{green!25} & \cellcolor{green!25} & \cellcolor{green!25} &                      \\
%             \cline{2-21}
%                                                                                 & Aplicación de técnicas estadísticas para la evaluación de precisión &                                       &                                       &                                      &                                       &                      &                      &                      &                      &                      &                      &                      &                      & \cellcolor{green!25} & \cellcolor{green!25} &                      \\
%             \cline{2-21}
%                                                                                 & Análisis de resultados y determinación de fiabilidad de los datos   &                                       &                                       &                                      &                                       &                      &                      &                      &                      &                      &                      &                      &                      & \cellcolor{green!25} & \cellcolor{green!25} &                      \\
%             \hline
%             \multirow{4}{0.41\textwidth}{Desarrollar  un sistema web y móvil para la obtención de datos
%                 georreferenciados de delitos con los cuales desarrollar un modelo analítico de
%                 Business Intelligence (BI) orientado al seguimiento de las tendencias de siniestros
%             en la ciudad de Ambato.}                                            & Desarrollo backend                                                  &                                       & \cellcolor{green!25}                  & \cellcolor{green!25}                 & \cellcolor{green!25}                  & \cellcolor{green!25} & \cellcolor{green!25} & \cellcolor{green!25} & \cellcolor{green!25} & \cellcolor{green!25} & \cellcolor{green!25} & \cellcolor{green!25} & \cellcolor{green!25} &                      &                      &                      \\
%             \cline{2-21}
%                                                                                 & Diseño del sistema web                                              &                                       & \cellcolor{green!25}                  & \cellcolor{green!25}                 & \cellcolor{green!25}                  & \cellcolor{green!25} & \cellcolor{green!25} & \cellcolor{green!25} & \cellcolor{green!25} & \cellcolor{green!25} & \cellcolor{green!25} & \cellcolor{green!25} & \cellcolor{green!25} &                      &                      &                      \\
%             \cline{2-21}
%                                                                                 & Diseño del sistema móvil                                            &                                       & \cellcolor{green!25}                  & \cellcolor{green!25}                 & \cellcolor{green!25}                  & \cellcolor{green!25} & \cellcolor{green!25} & \cellcolor{green!25} & \cellcolor{green!25} & \cellcolor{green!25} & \cellcolor{green!25} & \cellcolor{green!25} & \cellcolor{green!25} &                      &                      &                      \\
%             \cline{2-21}
%                                                                                 & Desarrollo del sistema web                                          &                                       & \cellcolor{green!25}                  & \cellcolor{green!25}                 & \cellcolor{green!25}                  & \cellcolor{green!25} & \cellcolor{green!25} & \cellcolor{green!25} & \cellcolor{green!25} & \cellcolor{green!25} & \cellcolor{green!25} & \cellcolor{green!25} & \cellcolor{green!25} &                      &                      &                      \\
%             \cline{2-21}
%                                                                                 & Desarrollo del sistema móvil                                        &                                       & \cellcolor{green!25}                  & \cellcolor{green!25}                 & \cellcolor{green!25}                  & \cellcolor{green!25} & \cellcolor{green!25} & \cellcolor{green!25} & \cellcolor{green!25} & \cellcolor{green!25} & \cellcolor{green!25} & \cellcolor{green!25} & \cellcolor{green!25} &                      &                      &                      \\
%             \cline{2-21}
%                                                                                 & Implementación del modelo de Business Intelligence                  &                                       &                                       &                                      &                                       &                      &                      &                      &                      &                      &                      &                      &                      & \cellcolor{green!25} & \cellcolor{green!25} &                      \\
%             \cline{2-21}
%                                                                                 & Pruebas y validación del sistema                                    &                                       &                                       &                                      &                                       &                      &                      &                      &                      &                      &                      &                      &                      & \cellcolor{green!25} & \cellcolor{green!25} &                      \\
%             \cline{2-21}
%                                                                                 & Presentación del trabajo final                                      &                                       &                                       &                                      &                                       &                      &                      &                      &                      &                      &                      &                      &                      &                      &                      & \cellcolor{green!25} \\
%             \hline
%             \insertTableNotes
%         \end{longtable}
%     \end{ThreePartTable}
% \end{landscape}


\newpage
\section{Bibliografía}
\bibliographystyle{biblio/IEEEtran}
\bibliography{biblio/main}


% \vfill
% \begin{center}
% 	\begin{tabular}{l}
% 		Revisado por:               \\
% 		\vspace{3cm}                \\
% 		\rule{0.5\textwidth}{0.5pt} \\
% 		Ing. Santiago Jara          \\
% 	\end{tabular}
% \end{center}


\end{document}
11 de marzo